\documentclass[12pt,a4paper]{article}
\usepackage{fontspec}
\usepackage{polyglossia}
\usepackage{graphicx}
\usepackage{hyperref}
\usepackage{geometry}
\geometry{a4paper, margin=2.5cm}

% Set up languages
\setdefaultlanguage{arabic}
\setotherlanguage{english}

% Font setup with Vazir
\newfontfamily\arabicfont[Script=Arabic]{Vazir}
\newfontfamily\englishfont{Times New Roman}
\setmainfont{Times New Roman}

\title{مصاحبهٔ طراحی سیستم: راهنمای یک فرد داخل‌کاری}
\author{الکس شو}
\date{}

\begin{document}

\maketitle

\begin{Arabic}

\section*{اطلاعات کتاب}

\textbf{مصاحبهٔ طراحی سیستم: راهنمای یک فرد داخل‌کاری}

تمام حقوق محفوظ است. این کتاب یا هر بخشی از آن بدون اجازهٔ کتبی ناشر، به هیچ شکلی قابل تکثیر یا استفاده نیست مگر برای نقل‌قول‌های کوتاه در نقد و بررسی کتاب.

\section*{دربارهٔ نویسنده}

الکس شو یک مهندس نرم‌افزار و کارآفرین باتجربه است. او پیش‌تر در توییتر، اپل، زیانگا و اوراکل کار کرده است. او مدرک کارشناسی ارشد خود را از دانشگاه کارنگی ملون دریافت کرده و علاقه‌مند به طراحی و پیاده‌سازی سیستم‌های پیچیده است.

\section*{اطلاعات تماس}

در صورت تمایل به اطلاع‌رسانی دربارهٔ انتشار فصل‌های جدید، لطفاً در لیست ایمیل ما عضو شوید:

\end{Arabic}

\url{https://bit.ly/3dtIcsE}

\begin{Arabic}

برای اطلاعات بیشتر، با آدرس زیر تماس بگیرید:

\end{Arabic}

\href{mailto:systemdesigninsider@gmail.com}{systemdesigninsider@gmail.com}

\begin{Arabic}

\textit{ویراستار:} پل سلیمان

\section*{نکات ترجمه}

\begin{itemize}
\item System Design Interview به «مصاحبهٔ طراحی سیستم» ترجمه شد.
\item An Insider's Guide به «راهنمای یک فرد داخل‌کاری» ترجمه شد تا مفهوم تخصصی بودن محتوا منتقل شود.
\item نام شرکت‌ها (Twitter, Apple, etc.) و دانشگاه (Carnegie Mellon) به شکل اصلی باقی ماندند.
\item M.S. به «کارشناسی ارشد» ترجمه شد.
\end{itemize}

\end{Arabic}

\end{document}
